\section{Characteristic two: the ordered Hessian}\label{sec:hessian}

In characteristic two the divided binary-cubic discriminant becomes
the doubled quadric $(AD+BC)^2$, so its pullback cannot distinguish
generic splitting.  The correct replacement retains an ordered
residual root.

Let $\ell=\PP\langle U,V\rangle\subset\PP(\Gamma^3E)$ and write
\[
\Hess^{\mathrm{div}}(uU+vV)
 =N_u(u,v)X^2+N_s(u,v)XY+D(u,v)Y^2.
\]
For
\[
 U=(A,B,C,D_3),\qquad V=(a,b,c,d),
\]
the three coefficient quadratics are
\[
\begin{aligned}
N_u&=(AC+B^2)u^2+(Ac+aC)uv+(ac+b^2)v^2,\\
N_s&=(AD_3+BC)u^2
 +(Ad+aD_3+Bc+bC)uv+(ad+bc)v^2,\\[-6pt]
D&=(BD_3+C^2)u^2+(Bd+bD_3)uv+(bd+c^2)v^2.
\end{aligned}
\]
The ordered-Hessian incidence
\begin{equation}\label{eq:hessian}
N_u(u,v)X^2+N_s(u,v)XY+D(u,v)Y^2=0
\end{equation}
has bidegree $(2,2)$ and arithmetic genus one.  Off $N_s=0$, its
Artin--Schreier class is $DN_u/N_s^2$.

The parameter space of cubic lines is
$G(1,\Gamma^3E)$.  Let
\[
\Sigma=\{\PP(QE):[Q]\in\PP(\Sym^2E)\}\subset G(1,\Gamma^3E)
\]
be the Veronese surface of common-quadratic lines.  In the ambient
cubic space, let $\mathcal Q_0=\{AD+BC=0\}\subset\PP(\Gamma^3E)$.
Its two families of lines give two Fano ruling conics in the
Grassmannian; one is the tangent-conic boundary of $\Sigma$, and we
call the other $\cN$.  A root-compatible pullback is the family of
lines obtained by fixing a split factor $R$ and varying the final two
linear factors in the consecutive Hankel contraction of a syndrome.

\begin{figure}[H]
\centering
\[
\begin{array}{ccccc}
G(1,\Gamma^3E)
&\supset&
\text{degenerate ordered-Hessian locus}
&=&
\Sigma\ \cup\ \operatorname{Fano}_1(\mathcal Q_0)\\[4pt]
&&\downarrow\ \text{root-compatible pullback}&&\\[4pt]
&&
\begin{array}{c}
\text{persistent component}\\
\cup\\
\text{Lucas-nucleus component}
\end{array}
&&
\cN:\ \text{no rank-two pullback}.
\end{array}
\]
\caption{Ordered-Hessian geometry.  The degeneracy locus is computed
in the Grassmannian first; root compatibility then selects the
contained syndrome carriers and eliminates the complementary ruling.}
\label{fig:hessian-flow}
\end{figure}

\begin{theorem}[Ordered-Hessian geometric degeneracy]
\label{thm:hessian}
After removing the forced vertical factors from intersections with the
twisted cubic, the complete geometric degeneracy locus of
\eqref{eq:hessian} consists of:
\begin{enumerate}[label=\textup{(\roman*)}]
\item the Veronese surface $\Sigma$ of lines $\PP(QE)$, on which the incidence is
      separably reducible; and
\item the two Fano ruling conics of $\mathcal Q_0$, on which the root
      projection of \eqref{eq:hessian} is inseparable.
\end{enumerate}
The complementary ruling $\cN$ has no nontrivial rank-two root-compatible
Hankel pullback.
\end{theorem}

\begin{proof}
A vertical factor divides $N_u,N_s,D$ exactly when the divided Hessian
vanishes at the corresponding point of $\ell$.  Its base scheme is the
twisted cubic, so the vertical factors are precisely
$\ell\cap C_3$.  A horizontal factor fixes one Hessian root for every
cubic on $\ell$.  The cubics with that root form the rank-three quadric
cone swept by secants through the corresponding point of $C_3$; every
line on the cone is a generator through its vertex.  Thus a horizontal
factor is already a chord/tangent line in $\Sigma$, or its collision
boundary.

After removing those factors, bidegree additivity forces any separable
factorization to be $(1,1)+(1,1)$.  Both factors are graphs of
projectivities.  Make the first graph the identity by independent
changes of coordinate on the parameter and root lines.  For a
semisimple second graph $g(t)=ct$, normalize the two cubics at its fixed
parameters to
\[
             U=(a,1,0,0),\qquad V=(0,0,\gamma,d).
\]
Comparison with
\[
 (uX+vY)(uX+cvY)
 =u^2X^2+(1+c)uvXY+cv^2Y^2
\]
gives
\[
 d=0,\qquad a=0,\qquad\gamma=1+c,\qquad\gamma^2=c.
\]
Hence $c^2+c+1=0$, and the normalized line is
$\PP\langle X^2Y,XY^2\rangle=XY\cdot\PP(E)$.  Undoing the
normalization gives exactly $\PP(QE)$.  In the unipotent normal form
$g(t)=t+1$, comparison with
\[
             (uX+vY)((u+v)X+vY)
\]
instead gives $d=\gamma=0$ and $a\gamma=1$, a contradiction.  The
identity case is included in the semisimple comparison and fails the
same equations.  Thus $\Sigma$ is the complete separable external
locus.

For the inseparable locus, $N_s$ vanishes identically exactly when
$\ell$ is a line on the smooth quadric
$\mathcal Q_0=\{AD+BC=0\}$.  Its Fano scheme has the two stated
ruling conics.  This classifies inseparable root projection, whether
or not the total $(2,2)$ curve is reduced.  The exact Pl\"ucker
equations and the constrained calculation excluding the complementary
ruling are given in Proposition~\ref{prop:hessian-pullback}.
\end{proof}

\begin{proposition}[Root-compatible pullback]\label{prop:hessian-pullback}
On the root-compatible Hankel parameter space, the reduced pullback of
$\Sigma$ is exactly the persistent catalecticant/Lucas-nucleus union.
The complementary ruling $\cN$ has no nontrivial rank-two pullback.
Intersections with the twisted cubic contribute only the vertical
factors removed before the reduced incidence is formed.
\end{proposition}

\begin{proof}
Use Pl\"ucker coordinates
\[
 (z_0,\ldots,z_5)=(p_{01},p_{02},p_{03},p_{12},p_{13},p_{23}).
\]
For $Q=aX^2+bXY+cY^2$, the line $QE$ has coordinates
\[
 [z_0:\cdots:z_5]
 =[a^2:ab:ac:b^2+ac:bc:c^2].
\]
Consequently $\Sigma$ is cut out scheme-theoretically by the
$2\times2$ minors of
\[
\begin{pmatrix}
z_0&z_1&z_2\\
z_1&z_3+z_2&z_4\\
z_2&z_4&z_5
\end{pmatrix}.                                             \tag{38}
\]
The tangent ruling conic is
\[
 z_1=z_4=0,\quad z_2=z_3,\quad z_0z_5=z_2^2,
\]
the square-quadratic boundary of the displayed parametrization.  The
complementary ruling is
\[
 \cN:\quad z_0=z_5=0,\quad z_2=z_3,\quad z_1z_4=z_2^2.      \tag{39}
\]
Together with the Chow equation for lines meeting $C_3$, (38)--(39)
are the complete universal degeneracy ideal found in the theorem.

For a split base $R$, the root-compatible line
$\ell_{f,R}=L_f(\PP(RE^\vee))$ has Pl\"ucker coordinates quadratic
in the coefficients of $R$.  If its generic point lies in $\Sigma$,
the minors (38) say that every iterated contraction pencil has a common
quadratic.  Proposition~\ref{prop:contained-rank-two}, applied to the
overlapping contraction flags, gives the persistent rank-two carrier;
when a consecutive contraction coefficient vanishes in the ground
characteristic, the same overlap equations and Lucas' theorem give
exactly the contraction-kernel nucleus flags.  Conversely, the
catalecticant and Lucas kernels make every minor in (38) vanish.
This proves equality of the reduced pullback with the stated union,
not only set-theoretic containment in one direction.

For $\cN$, the linear equations in (39), after substitution, give
the consecutive identities
\[
h_{-1}(R)^2=h_{-2}(R)h_0(R),\qquad
h_1(R)^2=h_0(R)h_2(R).
\]
If these hold for every split $R$, they hold polynomially because
split forms are Zariski dense over the algebraic closure.  The
contraction forms $h_i$ are linear forms in the coefficients of $R$.
In their UFD, an identity $L_1^2=L_0L_2$ implies either a zero form or
proportionality of the three nonzero forms: every irreducible factor
of $L_0$ or $L_2$ must be associated to $L_1$.  Applying this to the
two overlapping triples propagates proportionality across
$h_{-2},\ldots,h_2$.  Zero-form cases only lower the same contraction
rank.  Hence $\operatorname{rank}\ell_{f,R}\leq1$, so $\cN$ has no
nontrivial rank-two pullback.

Finally, the divided Hessian vanishes exactly on $C_3$.
A line--twisted-cubic intersection is therefore exactly the common
vertical factor removed above.  Scheme-theoretic division by that
factor commutes with field extension and cannot create a reduced
moving component.
\end{proof}

\begin{proposition}[Effective Hessian corollary]
\label{prop:hessian-effective}
For syndrome degree $n\geq5$, outside the persistent/Lucas union an
integral reduced slice exists whenever
\[
q\geq
\min\left\{\frac{(n-4)(n+11)}2+1,\ 9(n-4)\right\}.
\]
Its deletion degree is at most $3n-4$.  If also
$q+1-2\sqrt q>3n-4$, every split-free syndrome lies in that union.
\end{proposition}

\begin{proof}
For a split base $R$ of degree $n-4$, the root-compatible Pl\"ucker
coordinates are quadratic in its coefficients.  Outside the contained
union, choose one nonzero pulled-back equation $A$ of $\Sigma$ and one
nonzero pulled-back equation $B$ of $\cN$.  Each has individual
root degree at most four, and the single product $AB$ vanishes on the
whole reduced bad union and has individual root degree at most eight.
Vertical factors have already been removed by
Proposition~\ref{prop:hessian-pullback}.

Put $m=n-4$.  Schwartz--Zippel applied to $AB$ after adjoining the
Vandermonde requires
\[
q>8m+\binom m2=\frac{m(m+15)}2,
\]
which gives the first bound.  Alternatively partition $9m$ field
elements into disjoint nine-element blocks.  Successive univariate
interpolation shows that a nonzero polynomial of individual degree at
most eight cannot vanish on their product, giving the second bound.
The moving diagonal, determinant, branch, fixed-root collision, and
final collision divisors have degrees
\[
       n-4,\ 2,\ 2,\ 2(n-4),\ 4,
\]
whose sum is $3n-4$.  The normalized integral $(2,2)$ slice has genus at
most one, so the point bound leaves a rational ordered residual root
outside every deletion, and contraction lifts it squarefreely.
\end{proof}

Theorem~\ref{thm:hessian} is a containment theorem.  It does not assert
that every Lucas-carrier point is deep; that remains an arithmetic
question for each level.
